\chapter{न्यूटन के नियम}\label{ch:g12:newtons-laws}

लिफ़्ट में चोरी-छिपे एक तराज़ू ले जाइए। जैसे ही केबिन ऊपर की ओर खिंचता है, सुई
तीन \omterm{def:g10:orders-of-magnitude:unit}{किलोग्राम} भारी बताती है; ऊपर पहुँचते-पहुँचते हलका; और दोनों के बीच,
साफ़ सच --- जबकि उसमें बदला कुछ भी नहीं, केवल गति। यह अध्याय उस सुई के पीछे के
तीनों नियम रखता है और उन्हें उस विधि में ढाल देता है जो पूरी यांत्रिकी चलाती
है: निकाय चुनिए, उसके बल खींचिए, $\sum \vect F = m\vect a$ का प्रक्षेप कीजिए।

\section{संवेग, और पहला नियम ठीक-ठीक}

\begin{definition}[संवेग]\label{def:g12:newtons-laws:momentum}
वेग $\vect v$ से चलते द्रव्यमान $m$ (\unit{kg}) के पिंड का
\emph{संवेग}\index{संवेग} वह सदिश है
\[
\vect p = m \vect v \qquad (\unit{kg.m/s}):
\]
यानी द्रव्यमान से तौला हुआ वेग; और किसी निकाय के लिए, उसके भागों पर सदिश योग।
\end{definition}

\begin{definition}[जड़त्वीय निर्देश तंत्र]\label{def:g12:newtons-laws:inertial}
\emph{जड़त्वीय निर्देश तंत्र}\index{जड़त्वीय निर्देश तंत्र} वह है जिसमें हर
विलगित पिंड --- जिस पर कोई बल न हो, या जिसके बल संतुलित
हों --- सरल रेखा में अचर वेग से चले (\cref{ch:g10:inertia})। ज़मीन बहुत अच्छी परिशुद्धता तक
ऐसी ही है, और उस पर सरल रेखा में \omterm{def:g11:electric-gravitational-fields:capacitor}{एकसमान} चलता कोई भी केबिन भी; पर ब्रेक लगाती
रेलगाड़ी या घूमता हिंडोला नहीं: वहाँ खुला रखा सामान बिना धकेले ही त्वरित हो
जाता है।
\end{definition}

\begin{theorem}[न्यूटन का पहला नियम]\label{thm:g12:newtons-laws:first}\index{न्यूटन का पहला नियम}
जड़त्वीय तंत्र में किसी विलगित निकाय का \omterm{def:g12:newtons-laws:momentum}{संवेग} नहीं बदलता: $\vect p$ अचर
रहता है, यानी $d\vect p/dt = \vect 0$।
\end{theorem}
\admitted

\begin{remark}[जड़त्व, दो बार धार दिया हुआ]\label{rem:g12:newtons-laws:sharpened}
\cref{ch:g10:inertia} का \omterm{thm:g10:inertia:principle}{जड़त्व का सिद्धांत}, दो बार उन्नत: अब वह बताता है कि वह \emph{कहाँ}
लागू है --- वह जड़त्वीय तंत्रों की परिभाषा ही उन्हीं तंत्रों के रूप में कर देता
है जहाँ वह लागू है --- और वह एक सदिश, $m\vect v$, की रखवाली करता है, जिसका बहीखाता
अध्याय के अंत में रंग लाता है।
\end{remark}

\section{न्यूटन का दूसरा नियम}

\begin{theorem}[न्यूटन का दूसरा नियम]\label{thm:g12:newtons-laws:second}\index{न्यूटन का दूसरा नियम}
जड़त्वीय तंत्र में किसी पिंड पर लगता \omterm{def:g11:forces-and-motion:netforce}{परिणामी बल} उसके \omterm{def:g12:newtons-laws:momentum}{संवेग} के
बदलने की दर के बराबर होता है; और अचर द्रव्यमान के लिए, चूँकि $\vect p = m\vect v$,
\[
\sum \vect F = \frac{d\vect p}{dt} = m\,\frac{d\vect v}{dt} = m \vect a .
\]
\end{theorem}
\admitted

\begin{remark}[प्रकृति का एक नियम, और उसका मात्रक]\label{rem:g12:newtons-laws:nature}
इसकी कोई उपपत्ति नहीं है: यह प्रकृति का नियम है, जिसे तीन सदियों से प्रयोगों
के तराज़ू पर तौला जाता रहा है। पिछले वर्ष (\cref{ch:g11:forces-and-motion}) ने उसकी छाया दी थी, $\Delta v \propto F\,\Delta t/m$;
और अवकलज (\cref{ch:g12:kinematics-2d}) उसे हर क्षण के लिए ठीक-ठीक बना देता है --- उसकी बारीक शर्तें
वर्ष 1 के खंड में प्रतीक्षा कर रही हैं। बल को एक मात्रक मिल
जाता है, $\unit{N} = \unit{kg.m/s^2}$; और केवल अपने भार के नीचे पड़ा पिंड $m\vect a = m\vect g$ मानता है:
हर द्रव्यमान $\vect a = \vect g$ से गिरता है, और $g = \qty{9.81}{N/kg} = \qty{9.81}{m/s^2}$ --- दोनों पाठ्यांक एक ही हैं।
\end{remark}

\begin{example}[औसत बल, ब्योरे की ज़रूरत नहीं]\label{ex:g12:newtons-laws:average}
\qty{1300}{kg} की कार \omterm{def:g10:relative-motion:frame}{विराम} से \qty{8.0}{s} में \qty{100}{km/h} (\qty{27.8}{m/s}) तक पहुँच जाती है: बीच का
पल-पल का क़िस्सा चाहे जो हो, औसत \omterm{def:g11:forces-and-motion:netforce}{परिणामी बल} $\Delta p/\Delta t = 1300 \times 27.8/8.0 \approx \qty{4.5e3}{N}$ है --- \omterm{def:g12:newtons-laws:momentum}{संवेग}
किसी बल के प्रश्न को पहले-और-बाद के बहीखाते में बदल देता है।
\end{example}

\section{न्यूटन का तीसरा नियम}

\begin{theorem}[न्यूटन का तीसरा नियम]\label{thm:g12:newtons-laws:third}\index{न्यूटन का तीसरा नियम}
यदि पिंड $A$ पिंड $B$ पर बल $\vect F_{A \to B}$ लगाए, तो $B$ भी $A$ पर
बल $\vect F_{B \to A} = -\vect F_{A \to B}$ लगाता है: वही क्रिया-रेखा, वही परिमाण, उलटी दिशाएँ --- और वह
भी हर क्षण, गति चाहे जो हो, संपर्क हो या दूरी पर क्रिया।
\end{theorem}
\admitted

\begin{example}[किताब, मेज़ और पृथ्वी]\label{ex:g12:newtons-laws:book}
मेज़ पर रखी किताब: दो अंतःक्रियाएँ, दो युग्म। गुरुत्वीय: पृथ्वी किताब को नीचे
खींचती है ($\vect P$), किताब पृथ्वी को ऊपर। संपर्क: मेज़ किताब को ऊपर धकेलती है
($\vect N$), किताब मेज़ को नीचे दबाती है। चलना भी वही सौदा है: ज़मीन को पीछे
धकेलिए; और उसकी प्रतिक्रिया ही वह अकेला बल है जो आपको आगे ले जाता है।
\end{example}

\begin{omfigure}
\begin{tikzpicture}[scale=0.9, font=\small]
  \draw[thick] (0,0) -- (4.6,0); \draw (0.5,0) -- (0.5,-1.4) (4.1,0) -- (4.1,-1.4);
  \draw[fill=omExo!15] (1.6,0) rectangle (2.6,0.6);
  \draw[-Stealth, omThm, very thick] (1.9,0.3) -- ++(0,-1.7) node[below] {$\vect P$};
  \draw[-Stealth, omProp, very thick] (2.3,0.3) -- ++(0,1.7) node[above] {$\vect N$};
  \draw[-Stealth, omProp, very thick] (3.4,-0.05) -- ++(0,-1.7) node[below] {$-\vect N$};
  \draw[thick] (0,-3.1) -- (4.6,-3.1);
  \foreach \x in {0.4,0.8,...,4.5} \draw[gray] (\x,-3.1) -- (\x-0.25,-3.4);
  \node[above right] at (3.9,-3.05) {पृथ्वी};
  \draw[-Stealth, omThm, very thick] (0.9,-3.1) -- ++(0,1.7) node[above] {$-\vect P$};
\end{tikzpicture}

{\small एक किताब, दो अंतःक्रियाएँ, चार तीर: गुरुत्वीय युग्म (पृथ्वी किताब को
खींचती है, किताब पृथ्वी को) और संपर्क-युग्म (मेज़ किताब को धकेलती है, किताब
मेज़ को दबाती है)। $\vect P$ और $\vect N$ साथी \emph{नहीं} हैं: वे एक ही किताब पर
लगते हैं --- और केवल इसलिए कट जाते हैं कि उसका त्वरण शून्य है।}
\end{omfigure}

\section{विधि: चुनिए, खींचिए, प्रक्षेप कीजिए}

\begin{method}[गतिकी की समस्या हल करना]\label{met:g12:newtons-laws:method}
\begin{enumerate}
  \item \emph{निकाय}: बताइए कि किस पिंड (या पिंडों के समुच्चय) का अध्ययन हो
        रहा है।
  \item \emph{तंत्र}: कोई जड़त्वीय तंत्र चुनिए और गति के अनुरूप अक्ष ---
        ढलान वाले अक्ष, या वक्र के लिए \omterm{def:g12:kinematics-2d:frenet}{फ्रेने आधार} (\cref{ch:g12:kinematics-2d})।
  \item \emph{सूची}: हर बल गिनाइए --- पहले भार, फिर हर संपर्क
        के लिए एक बल, और कुछ नहीं --- और मुक्त-पिंड आरेख खींचिए: एक
        बिंदु, और हर बल के लिए एक तीर।
  \item \emph{प्रक्षेप}: $\sum \vect F = m\vect a$ लिखिए और उसे हर अक्ष पर प्रक्षेपित कीजिए।
  \item \emph{हल और जाँच}: मात्रक, चिह्न, और सीमांत स्थितियाँ
        ($\mu \to 0$, $\alpha \to 0$, \dots)।
\end{enumerate}
\end{method}

\begin{definition}[अभिलंब बल और घर्षण]\label{def:g12:newtons-laws:friction}
कोई सतह किसी पिंड पर दो घटकों से लगती है: सतह के लंबवत \emph{अभिलंब
बल}\index{अभिलंब बल} $\vect N$, और स्पर्शरेखीय \omterm{def:g10:inertia:inventory}{घर्षण} $\vect f$। प्रयोग
\omterm{def:g10:inertia:inventory}{घर्षण} के नियम देते हैं: जब तक पिंड सरकता नहीं, \omterm{def:g11:forces-and-motion:friction}{स्थैतिक घर्षण}
एक छत तक अपने को बढ़ाता जाता है, $f_s \leq \mu_s N$; और सरकना शुरू होते ही \omterm{def:g11:forces-and-motion:friction}{गतिज
घर्षण} $f_k = \mu_k N$ होता है, सरकने के विरुद्ध। \emph{\omterm{def:g10:inertia:inventory}{घर्षण} के
गुणांक}% \index{घर्षण गुणांक} $\mu_s$ और $\mu_k$ ($\mu_k \leq \mu_s$) केवल दोनों
पदार्थों पर निर्भर करते हैं।
\end{definition}

\begin{example}[ढलान, दो बार]\label{ex:g12:newtons-laws:incline}
कोई गुटका कोण $\alpha$ की ढलान पर नीचे सरकता है; अक्ष: $x$ ढलान के नीचे की
ओर, $y$ लंबवत। $y$ पर: $N = mg\cos\alpha$। घर्षणहीन ढलान पर $x$ के लिए: $ma = mg\sin\alpha$,
इसलिए $a = g\sin\alpha$ --- कोई द्रव्यमान नहीं; और $\alpha = 30^\circ$ पर $a = \qty{4.9}{m/s^2}$। \omterm{def:g11:forces-and-motion:friction}{गतिज घर्षण}
के साथ: $a = g(\sin\alpha - \mu_k\cos\alpha)$; और $\mu_k = 0.20$ के लिए $a = 9.81\,(0.500 - 0.173) = \qty{3.2}{m/s^2}$। और सबसे पहले वह टिका ही तब तक रहा
जब तक ज़रूरी \omterm{def:g11:forces-and-motion:friction}{स्थैतिक घर्षण} अपनी छत के नीचे अँटता रहा: $\tan\alpha \leq \mu_s$।
\end{example}

\begin{omfigure}
\begin{tikzpicture}[scale=1.0, font=\small]
  \draw[gray!60] (-0.3,0) -- (6.4,0); \draw[thick] (0,0) -- (30:6.2);
  \draw (1.6,0) arc (0:30:1.6); \node at (14:2.0) {$\alpha$};
  \begin{scope}[shift={(30:3.6)}, rotate=30]
    \draw[fill=omExo!15] (-0.6,0) rectangle (0.6,0.75);
    \coordinate (G) at (0,0.375); \fill (G) circle (1.5pt);
    \draw[-Stealth, omDef, dashed, thick] (G) -- ++(-2.3,0) node[above] {$x$};
    \draw[-Stealth, omDef, dashed, thick] (G) -- ++(0,1.9) node[right] {$y$};
    \draw[-Stealth, omThm, very thick] (G) -- ++(0,1.5) node[left] {$\vect N$};
    \draw[-Stealth, omThm, very thick] (G) -- ++(1.35,0) node[above] {$\vect f$};
  \end{scope}
  \draw[-Stealth, omThm, very thick] (G) -- ++(0,-1.9) node[below] {$\vect P$};
  \draw[-Stealth, omExo, dashed] (G) -- ++(210:0.95) (G) -- ++(-60:1.645);
\end{tikzpicture}

{\small ढलान पर मुक्त-पिंड आरेख: भार $\vect P$, \omterm{def:g12:newtons-laws:friction}{अभिलंब बल}
$\vect N$, और नीचे सरकते गुटके के लिए ढलान पर ऊपर की ओर \omterm{def:g10:inertia:inventory}{घर्षण} $\vect f$।
ढलान के अनुदिश और लंबवत अक्ष $\vect P$ को $mg\sin\alpha$ और $mg\cos\alpha$ में बाँट देते हैं
(बिंदुदार)।}
\end{omfigure}

\begin{definition}[आभासी भार]\label{def:g12:newtons-laws:apparent}
किसी पिंड का \emph{आभासी भार}\index{आभासी भार} उसके और उसके आधार के बीच का
\omterm{def:g12:newtons-laws:friction}{अभिलंब बल} है --- यानी उसके नीचे रखा तराज़ू जो पढ़ता है।
\end{definition}

\begin{example}[लिफ़्ट]\label{ex:g12:newtons-laws:elevator}
\qty{70}{kg} की सवारी \omterm{def:g10:universal-gravitation:weight}{ऊर्ध्वाधर} त्वरण $a_z$ (ऊपर की ओर धनात्मक) वाली लिफ़्ट में
तराज़ू पर खड़ी है: $N - mg = ma_z$, इसलिए $N = m(g + a_z)$। $a_z = \qty{1.5}{m/s^2}$ से ऊपर खिंचते समय: \qty{792}{N},
भारी; ऊपर पहुँचते-पहुँचते ब्रेक लगने पर: \qty{582}{N}, हलकी; और एक-सी चाल पर
$mg = \qty{687}{N}$ --- यानी पहला नियम; और मुक्त पतन में, $a_z = -g$: $N = 0$, यानी तराज़ू पर
भारहीनता।
\end{example}

\begin{example}[दो गुटके और एक घिरनी]\label{ex:g12:newtons-laws:pulley}
घर्षणहीन मेज़ पर रखा गुटका $m_1 = \qty{4.0}{kg}$ एक हलकी रस्सी से, आदर्श घिरनी के ऊपर से
होकर, लटकते गुटके $m_2 = \qty{1.0}{kg}$ से बँधा है। रस्सी दोनों सिरों पर वही \omterm{def:g10:inertia:inventory}{तनाव}
$T$ पहुँचाती है; और दोनों गुटकों के त्वरण का परिमाण एक ही रहता है $a$।
हर एक के लिए, उसकी अपनी गति के अनुदिश एक-एक दूसरा नियम: $m_1 a = T$ और $m_2 a = m_2 g - T$,
इसलिए $a = m_2\,g/(m_1 + m_2) = \qty{1.96}{m/s^2}$ और $T = m_1 a = \qty{7.8}{N}$ --- जो $m_2 g = \qty{9.8}{N}$ से छोटा है, और $m_2$ के नीचे की ओर
त्वरित होने के लिए ऐसा होना ही चाहिए।
\end{example}

\begin{omfigure}
\begin{tikzpicture}[scale=0.95, font=\small]
  \draw[thick] (0,1.1) -- (4.0,1.1); \draw (0.4,1.1) -- (0.4,-0.4) (3.6,1.1) -- (3.6,-0.4);
  \draw[fill=omExo!15] (1.0,1.1) rectangle (2.0,1.75); \node at (1.5,1.42) {$m_1$};
  \draw (4.0,1.1) -- (4.25,1.35);
  \draw[fill=white] (4.25,1.35) circle (0.25); \fill (4.25,1.35) circle (1.2pt);
  \draw[thick] (2.0,1.6) -- (4.25,1.6) (4.5,1.35) -- (4.5,0.15);
  \draw[fill=omExo!15] (4.2,-0.45) rectangle (4.8,0.15); \node at (4.5,-0.15) {$m_2$};
  \draw[-Stealth, omThm, very thick] (2.0,1.42) -- ++(1.2,0) node[pos=0.8, below] {$\vect T$};
  \draw[-Stealth, omThm, very thick] (4.5,0.15) -- ++(0,0.85) node[right] {$\vect T$};
  \draw[-Stealth, omThm, very thick] (4.5,-0.15) -- ++(0,-1.1) node[right] {$\vect P_2$};
  \draw[-Stealth, omDef, thick] (1.1,2.15) -- ++(0.8,0) node[right] {$\vect a$};
  \draw[-Stealth, omDef, thick] (5.4,0.4) -- ++(0,-0.8) node[below] {$\vect a$};
\end{tikzpicture}

{\small दो गुटके, एक हलकी रस्सी, एक आदर्श घिरनी: वही \omterm{def:g10:inertia:inventory}{तनाव} $T$
दोनों सिरों को खींचता है, और दोनों गुटकों के त्वरण का परिमाण एक ही रहता है ---
हर एक के लिए एक दूसरा नियम, यानी दो समीकरण, दो अज्ञात।}
\end{omfigure}

\begin{example}[समतल मोड़]\label{ex:g12:newtons-laws:curve}
द्रव्यमान $m$ की कार त्रिज्या $R$ के समतल मोड़ को अचर चाल $v$ से
घूमती है। \omterm{def:g12:kinematics-2d:frenet}{फ्रेने आधार} (\cref{ch:g12:kinematics-2d}) में त्वरण विशुद्ध \omterm{def:g10:refraction:angles}{अभिलंब} है, $a_N = v^2/R$,
और केंद्र की ओर ताका हुआ। \omterm{def:g10:universal-gravitation:weight}{ऊर्ध्वाधर} दिशा में $N = mg$; और क्षैतिज दिशा
में एकमात्र उपलब्ध अभिकेंद्र बल सड़क का टायरों पर लगता \omterm{def:g11:forces-and-motion:friction}{स्थैतिक
घर्षण} है: $f = mv^2/R \leq \mu_s N = \mu_s mg$, इसलिए $v \leq \sqrt{\mu_s\, g R}$, और वह भी द्रव्यमान चाहे जो हो।
$R = \qty{90}{m}$, सूखी सड़क ($\mu_s = 0.70$) के लिए: $v_{\max} \approx \qty{25}{m/s}$ (\qty{89}{km/h}); बर्फ़ पर वह छत ढह जाती
है --- \cref{exo:g12:newtons-laws:13} सड़क को झुका देता है; और \cref{exo:g12:newtons-laws:10} किसी तार पर वही प्रक्षेप झुलाता
है।
\end{example}

\begin{omfigure}
\begin{tikzpicture}[scale=0.85, font=\small]
  \draw[gray!50, dashed] (0,0) circle (2.1); \fill (0,0) circle (1.5pt);
  \draw[dashed, omExo] (0,0) -- (150:2.1) node[midway, above] {$R$};
  \draw[fill=omExo!15, rotate around={90:(2.1,0)}] (1.6,-0.28) rectangle (2.6,0.28);
  \draw[-Stealth, omDef, thick] (2.1,0.55) -- ++(0,1.4) node[above] {$\vect v$};
  \draw[-Stealth, omThm, very thick] (2.38,0) -- ++(-1.45,0) node[pos=0.85, above] {$\vect f$};
  \node at (0.2,-2.7) {ऊपर से दृश्य};
\end{tikzpicture}
\qquad
\begin{tikzpicture}[scale=0.85, font=\small]
  \draw[thick] (-0.5,0) -- (3.7,0); \draw[fill=omExo!15] (0.9,0.42) rectangle (2.5,1.25);
  \draw[fill=white] (1.2,0.42) circle (0.26) (2.2,0.42) circle (0.26);
  \coordinate (C) at (1.7,0.85); \fill (C) circle (1.2pt);
  \draw[-Stealth, omThm, very thick] (C) -- ++(0,1.4) node[above] {$\vect N$};
  \draw[-Stealth, omThm, very thick] (C) -- ++(0,-1.7) node[below] {$\vect P$};
  \draw[-Stealth, omThm, very thick] (C) -- ++(-1.6,0) node[above] {$\vect f$};
  \node at (1.6,-2.7) {पीछे से दृश्य (केंद्र बाईं ओर)};
\end{tikzpicture}

{\small अचर चाल से समतल मोड़ में कार: ऊपर से (बाएँ) वेग स्पर्शी है और
\omterm{def:g11:forces-and-motion:netforce}{परिणामी बल} केंद्र की ओर ताका है; पीछे से (दाएँ) $\vect N$ $\vect P$ को
संतुलित करता है, और सड़क का बग़ली \omterm{def:g11:forces-and-motion:friction}{स्थैतिक घर्षण} ही पूरा अभिकेंद्र
बल है।}
\end{omfigure}

\section{विलगित निकाय: संवेग संरक्षित रहता है}

\begin{proposition}[संवेग का संरक्षण]\label{prop:g12:newtons-laws:conservation}\index{संवेग संरक्षण}
जिन दो पिंडों पर बाहरी बल या तो लुप्त हों या संतुलित हों,
उनका कुल \omterm{def:g12:newtons-laws:momentum}{संवेग} $\vect p_1 + \vect p_2$ अचर रहता है, चाहे वे एक-दूसरे पर कैसे भी
बल लगाएँ।
\end{proposition}

\begin{proof}
दोनों दूसरे नियम जोड़ दीजिए: $\frac{d}{dt}(\vect p_1 + \vect p_2) = \vect F_{2 \to 1} +
\vect F_{1 \to 2} + \vect F_{\text{बाह्य}}$: भीतरी युग्म तीसरे नियम से कट जाता है, और
परिकल्पना से $\vect F_{\text{बाह्य}} = \vect 0$।
\end{proof}

\begin{example}[प्रतिक्षेप]\label{ex:g12:newtons-laws:recoil}
चिकनी बर्फ़ पर \omterm{def:g10:relative-motion:frame}{विराम} में खड़ी \qty{60}{kg} की स्केटर \qty{4.0}{kg} की एक भारी गेंद
क्षैतिज दिशा में \qty{6.0}{m/s} से फेंकती है। पहले: $\vect p = \vect 0$; बाद में: $0 = 60\,V + 4.0 \times 6.0$, इसलिए
$V = \qty{-0.40}{m/s}$ --- वह पीछे की ओर सरक जाती है। उसे बाहर से किसी ने नहीं धकेला: गेंद ने
धकेला (तीसरा नियम), और बहीखाता शून्य पर ही बना रहता है।
\end{example}

\begin{remark}[राकेट किस पर धक्का देता है]\label{rem:g12:newtons-laws:rocket}
राकेट द्रव्यमान फेंकने की मशीन है। उसके इंजन हर \omterm{def:g10:orders-of-magnitude:unit}{सेकंड} गैस की एक खेप पीछे की ओर
उछालते हैं; गैस पीछे की ओर \omterm{def:g12:newtons-laws:momentum}{संवेग} ले जाती है, इसलिए राकेट उतना ही आगे
की ओर पा लेता है --- यानी एक स्थिर \omterm{def:g10:inertia:inventory}{प्रणोद}, जो मोटे तौर पर प्रति \omterm{def:g10:orders-of-magnitude:unit}{सेकंड}
निकाले गए द्रव्यमान गुणा निकास-चाल है। राकेट अपने ही निकास पर धक्का देता है,
हवा पर नहीं: इसलिए वह ख़ाली अवकाश में सबसे अच्छा काम करता है।
\end{remark}

\section{अभ्यास}

\begin{exercise}[$\star$]\label{exo:g12:newtons-laws:1}
इनका \omterm{def:g12:newtons-laws:momentum}{संवेग} निकालिए: (क) \qty{10}{m/s} पर \qty{70}{kg} का धावक; (ख) \qty{50}{km/h} पर
\qty{1300}{kg} की कार; (ग) \qty{800}{m/s} पर \qty{8.0}{g} की गोली। इन्हें क्रम में रखिए। सबसे
तेज़ सबसे पीछे क्यों है?
\end{exercise}

\begin{exercise}[$\star$]\label{exo:g12:newtons-laws:2}
\qty{70}{kg} की सवारी लिफ़्ट में तराज़ू पर खड़ी है। तराज़ू क्या पढ़ेगा जब केबिन
(क) \qty{1.2}{m/s^2} से ऊपर की ओर त्वरित हो; (ख) अचर \qty{2.0}{m/s} से चढ़े; (ग) \qty{1.2}{m/s^2} से
नीचे की ओर त्वरित हो? (ख) का फ़ैसला कौन-सा नियम करता है?
\end{exercise}

\begin{exercise}[$\star$]\label{exo:g12:newtons-laws:3}
कोण $30^\circ$ की घर्षणहीन ढलान पर एक गुटका छोड़ा जाता है। उसका त्वरण निकालिए,
फिर \qty{3.0}{m} सरकने के बाद उसकी चाल। द्रव्यमान कहाँ चला गया?
\end{exercise}

\begin{exercise}[$\star$]\label{exo:g12:newtons-laws:4}
\qty{58}{g} की टेनिस गेंद \qty{15}{m/s} से आती है और उसी रेखा पर \qty{25}{m/s} से लौटा दी
जाती है; संपर्क \qty{5.0}{ms} रहता है। \omterm{def:g12:newtons-laws:momentum}{संवेग} का बदलाव और तारों का औसत
बल निकालिए; और उस बल की तुलना गेंद के भार से
कीजिए।
\end{exercise}

\begin{exercise}[$\star$]\label{exo:g12:newtons-laws:5}
मेज़ पर एक किताब रखी है। किताब पर लगते दोनों बल बताइए और हर एक का
क्रिया--प्रतिक्रिया साथी भी (किस पिंड पर? किस दिशा में?)। किताब का भार
और मेज़ का धक्का एक-दूसरे के साथी क्यों नहीं हैं?
\end{exercise}

\begin{exercise}[$\star\star$]\label{exo:g12:newtons-laws:6}
\qty{50}{kg} का संदूक क्षैतिज फ़र्श पर रखा है, $\mu_s = 0.45$, $\mu_k = 0.35$। (क) कितना क्षैतिज
धक्का उसे बस-बस चला देगा? (ख) वही धक्का बनाए रखने पर उसका त्वरण क्या होगा?
(ग) उसे अचर वेग से चलाए रखने के लिए कितना धक्का चाहिए?
\end{exercise}

\begin{exercise}[$\star\star$]\label{exo:g12:newtons-laws:7}
वही संदूक एक झुकाई जा सकने वाली ढलान पर रखा है ($\mu_s = 0.45$, $\mu_k = 0.35$)। (क) किस
कोण पर वह सरकना शुरू करता है? (ख) $30^\circ$ पर उसका त्वरण निकालिए। (ग) (क) वाले
कोण पर एक बार चल पड़ने के बाद क्या वह त्वरित होता ही रहता है? क्यों?
\end{exercise}

\begin{exercise}[$\star\star$]\label{exo:g12:newtons-laws:8}
घर्षणहीन मेज़ पर रखा \qty{3.0}{kg} का गुटका आदर्श घिरनी के ऊपर से होकर लटकते
\qty{2.0}{kg} के गुटके से बँधा है। त्वरण और \omterm{def:g10:inertia:inventory}{तनाव} निकालिए। \omterm{def:g10:inertia:inventory}{तनाव} को
लटकते भार से छोटा निकलना ही क्यों था?
\end{exercise}

\begin{exercise}[$\star\star$]\label{exo:g12:newtons-laws:9}
त्रिज्या \qty{50}{m} का मोड़ समतल है। $\mu_s = 0.70$ (सूखा) और $0.40$ (गीला) के लिए कार
की अधिकतम चाल \unit{km/h} में निकालिए। उत्तर स्कूटर और ट्रक, दोनों पर एक-सा क्यों
लागू होता है?
\end{exercise}

\begin{exercise}[$\star\star$]\label{exo:g12:newtons-laws:10}
\qty{1.2}{m} के तार पर लटका \qty{0.30}{kg} का गोलक क्षैतिज वृत्त खींचता है, और तार
\omterm{def:g10:universal-gravitation:weight}{ऊर्ध्वाधर} से अचर $25^\circ$ बनाए रखता है (शंक्वाकार लोलक)। दूसरे नियम का
\omterm{def:g10:universal-gravitation:weight}{ऊर्ध्वाधर} दिशा में प्रक्षेप कीजिए, फिर त्रिज्या के अनुदिश:
\omterm{def:g10:inertia:inventory}{तनाव} और चाल निकालिए।
\end{exercise}

\begin{exercise}[$\star\star$]\label{exo:g12:newtons-laws:11}
माल-लिफ़्ट में \qty{5.0}{kg} का पार्सल तराज़ू पर रखा है; पाठ्यांक ये हैं:
\begin{center}\small\begin{tabular}{lccc}
चरण & \qtyrange{0}{2}{s} & \qtyrange{2}{8}{s} & \qtyrange{8}{10}{s}\\
तराज़ू (\unit{N}) & 54.0 & 49.0 & 44.1
\end{tabular}\end{center}
हर चरण में त्वरण निकालिए और एक संभव यात्रा का वर्णन कीजिए। क्या लिफ़्ट पूरे
समय \emph{नीचे} की ओर चल रही हो सकती थी?
\end{exercise}

\begin{exercise}[$\star\star\star$]\label{exo:g12:newtons-laws:12}
चिकनी बर्फ़ पर \omterm{def:g10:relative-motion:frame}{विराम} में खड़ी \qty{60}{kg} की स्केटर \qty{4.0}{kg} की गेंद क्षैतिज दिशा
में फेंकती है। (क) गेंद \emph{ज़मीन के सापेक्ष} \qty{8.0}{m/s} से निकलती है: उसकी
प्रतिक्षेप-चाल? (ख) अब \qty{8.0}{m/s} \emph{उसके सापेक्ष} नापा गया है: दोनों चालें
फिर से निकालिए। (ग) उसे पीछे किस बल ने धकेला?
\end{exercise}

\begin{exercise}[$\star\star\star$]\label{exo:g12:newtons-laws:13}
त्रिज्या \qty{80}{m} वाला राजमार्ग का मोड़ \qty{20}{m/s} के लिए इस तरह बनाया गया है कि
\omterm{def:g10:inertia:inventory}{घर्षण} की कोई मदद न लगे। दिखाइए कि सड़क को कोण $\theta$ से झुका होना
चाहिए जो $\tan\theta = v^2/(gR)$ से मिलता है, और उसे निकालिए। जो कार मोड़ को उससे धीमे --- या
तेज़ --- लेती है, उसके लिए \omterm{def:g10:inertia:inventory}{घर्षण} को क्या करना पड़ता है?
\end{exercise}

\begin{exercise}[$\star\star\star$]\label{exo:g12:newtons-laws:14}
किसी राकेट के इंजन प्रति \omterm{def:g10:orders-of-magnitude:unit}{सेकंड} \qty{250}{kg} गैस \qty{2500}{m/s} से निकालते हैं। (क) प्रति
\omterm{def:g10:orders-of-magnitude:unit}{सेकंड} कितना \omterm{def:g12:newtons-laws:momentum}{संवेग} बाहर जाता है --- गैसें कितना \omterm{def:g10:inertia:inventory}{प्रणोद} लौटाती
हैं? (ख) प्रज्वलन के समय राकेट का द्रव्यमान \qty{5.0e4}{kg} है: उसका भार?
उसका उड़ान-भरने का त्वरण? (ग) यह इंजन ख़ाली अवकाश में भी उतना ही क्यों
काम करता है?
\end{exercise}

\begin{exercise}[$\star\star\star$]\label{exo:g12:newtons-laws:15}
मेज़ पर रखा \qty{4.0}{kg} का गुटका ($\mu_s = 0.35$, $\mu_k = 0.25$) आदर्श घिरनी के ऊपर से होकर
लटकते द्रव्यमान $m_2$ से बँधा है। (क) गति शुरू करने के लिए न्यूनतम $m_2$
कितना चाहिए? (ख) $m_2 = \qty{2.0}{kg}$ के लिए त्वरण और \omterm{def:g10:inertia:inventory}{तनाव} निकालिए। (ग) जाँचिए कि
$\mu_k = 0$ रखने पर \cref{exo:g12:newtons-laws:8} वाला घर्षणहीन सूत्र फिर मिल जाता है।
\end{exercise}

\section{समस्या: रज्जुमार्ग}

\begin{problem}[{सप्ताहांत समस्या --- रज्जुमार्ग का प्रमाणन: विराम में ढलान
जितना तनाव माँगती है, चलना शुरू करने का अधिभार, मीनार के ऊपर मिलती राहत,
आपात-रुकाव का दाम --- और वह एक संख्या जिसका वादा केबिन में लगी तख़्ती कर सकती
है}]
\label{pb:g12:newtons-laws:1}
पहाड़ का एक रज्जुमार्ग --- केबिन और गाड़ी मिलाकर ख़ाली द्रव्यमान $M_0 = \qty{2.0e3}{kg}$ ---
$\alpha = 30^\circ$ पर झुकी पटरी-रस्सी पर लुढ़कता है और ढलान के समांतर \omterm{def:g10:inertia:inventory}{तनाव}
$\vect T$ वाली खिंचाव-रस्सी से खींचा जाता है। लोटनी \omterm{def:g10:inertia:inventory}{घर्षण} नगण्य है;
केबिन का फ़र्श पूरे सफ़र में क्षैतिज बना रहता है। केबिन में औसतन \qty{80}{kg}
द्रव्यमान की $25$ सवारियाँ बैठती हैं; और प्रमाणित करने वाले अभियंता आप हैं।
जब तक कुछ और न कहा जाए, केबिन भरी हुई है: $M = \qty{4.0e3}{kg}$।

\medskip
\noindent\textbf{भाग I --- ढलान पर \omterm{def:g10:relative-motion:frame}{विराम} में।}
\begin{enumerate}
  \item निकाय चुनिए; बताइए कि पहाड़ का तंत्र जड़त्वीय तंत्र के रूप में क्यों
        चल जाएगा; और तीनों बल दिशाओं समेत गिनाइए।
  \item इस प्रतीक्षा पर कौन-सा नियम राज करता है? उसे ढलान के अनुदिश और लंबवत
        प्रक्षेपित करके $T$ और पटरी-रस्सी का \omterm{def:g12:newtons-laws:friction}{अभिलंब बल} $N$
        लिखिए।
  \item $T$ और $N$ निकालिए।
  \item ख़ाली केबिन के लिए $T$ फिर से निकालिए। $T$ कुल द्रव्यमान के
        समानुपाती क्यों है?
  \item खिंचाव-रस्सी टूट जाती है और ब्रेक भी नहीं हैं: दिखाइए कि भागती केबिन
        का त्वरण भरी हो या ख़ाली, $g \sin\alpha$ ही रहता है, और उसे निकालिए।
\end{enumerate}

\medskip
\noindent\textbf{भाग II --- चलना शुरू करना।} चरखी भरी केबिन को $a = \qty{0.60}{m/s^2}$ से
ढलान पर ऊपर त्वरित करती है, जब तक वह $v = \qty{6.0}{m/s}$ की एक-सी चाल न पकड़ ले।
\begin{enumerate}[resume]
  \item इस चरण में \omterm{def:g10:inertia:inventory}{तनाव} $T'$ लिखिए और निकालिए।
  \item चलना शुरू करना \omterm{def:g10:inertia:inventory}{तनाव} को \omterm{def:g10:relative-motion:frame}{विराम} वाले मान से कितने अंश ऊपर चढ़ा
        देता है?
  \item यह चरण कितनी देर और कितनी दूरी तक चलता है?
  \item \qty{80}{kg} की एक सवारी क्षैतिज फ़र्श पर खड़ी है। दूसरे नियम का
        \omterm{def:g10:universal-gravitation:weight}{ऊर्ध्वाधर} प्रक्षेप करके उसका \omterm{def:g12:newtons-laws:apparent}{आभासी भार} लिखिए और
        निकालिए; $mg$ से तुलना कीजिए।
  \item उसे क्षैतिज दिशा में कौन-सा बल त्वरित करता है, और कितना
        बड़ा? सवारियाँ झुक क्यों जाती हैं?
\end{enumerate}

\medskip
\noindent\textbf{भाग III --- मीनार के ऊपर।} अचर \qty{6.0}{m/s} पर पटरी-रस्सी एक
सहारा-मीनार के ऊपर त्रिज्या $R = \qty{40}{m}$ के \omterm{def:g10:universal-gravitation:weight}{ऊर्ध्वाधर} वृत्तीय चाप के अनुदिश मुड़
जाती है; और शिखर पर वेग क्षैतिज होता है।
\begin{enumerate}[resume]
  \item शिखर पर \omterm{def:g12:kinematics-2d:frenet}{फ्रेने आधार} (\cref{ch:g12:kinematics-2d}) में त्वरण के दोनों घटक उनके
        मानों समेत दीजिए।
  \item दूसरे नियम का \omterm{def:g10:refraction:angles}{अभिलंब} के अनुदिश प्रक्षेप करके गाड़ी पर रस्सी का
        \omterm{def:g12:newtons-laws:friction}{अभिलंब बल} $N'$ लिखिए और निकालिए।
  \item शिखर पर सवारी के \omterm{def:g12:newtons-laws:apparent}{आभासी भार} के लिए वही प्रश्न: वह कितने अंश
        हलकी हो जाती है?
  \item किस चाल पर $N'$ लुप्त हो जाता? उसकी तुलना यात्रा-चाल से कीजिए और
        टिप्पणी कीजिए।
  \item मीनार के ऊपर गाड़ी रस्सी पर कौन-सा बल लगाती है --- कौन-सा
        नियम यह कहता है, और उसका परिमाण तथा दिशा क्या है?
\end{enumerate}

\medskip
\noindent\textbf{भाग IV --- ब्रेक, एक छलाँग, और तख़्ती।}
\begin{enumerate}[resume]
  \item क्षैतिज आगमन-पट्टी पर एक आपात-रुकाव भरी केबिन को \qty{1.5}{s} में \qty{6.0}{m/s}
        से \omterm{def:g10:relative-motion:frame}{विराम} तक ले आता है: रुकाव भर के औसत के रूप में $\sum \vect F = d\vect p/dt$ पढ़कर
        ब्रेक-बल निकालिए।
  \item \qty{80}{kg} की सवारी को रोकने के लिए कितना क्षैतिज बल चाहिए, और
        क्या जूते--फ़र्श का \omterm{def:g11:forces-and-motion:friction}{स्थैतिक घर्षण} ($\mu_s = 0.40$) उसे दे सकता है?
        निष्कर्ष लिखिए।
  \item मंच पर, ब्रेक हटे हुए, ख़ाली केबिन मुक्त और निश्चल लटकी है; आख़िरी
        सवारी (\qty{80}{kg}) \emph{केबिन के सापेक्ष} \qty{2.5}{m/s} से क्षैतिज दिशा में
        कूद जाती है। केबिन की प्रतिक्षेप-चाल और सवारी की ज़मीन के सापेक्ष चाल
        निकालिए।
  \item एक वाक्य में: कौन-सा नोदन-यंत्र इसी प्रतिक्षेप-सिद्धांत पर जीता है,
        और उसका \omterm{def:g10:inertia:inventory}{प्रणोद} किससे तय होता है?
  \item \emph{प्रमाणन।} खिंचाव-रस्सी \qty{100}{kN} पर टूटती है; और नियम कार्यकारी
        \omterm{def:g10:inertia:inventory}{तनाव} को उसके पाँचवें हिस्से तक सीमित रखते हैं। दिखाइए कि ऊपर
        मिली सबसे बुरी स्थिति चलना शुरू करने वाली है, उससे अधिकतम कुल
        द्रव्यमान निकालिए, और तख़्ती छाप दीजिए: \qty{80}{kg} की कितनी सवारियाँ बैठ
        सकती हैं? सुरक्षा-गुणक जाँच लीजिए।
\end{enumerate}
\end{problem}
